\subsection{Representing Sequences}
\label{Section 2.2.1}

One of the more useful structures for use in procedures is a
\newterm{sequence}---an ordered collection of data objects.  In Scheme, this is
done with pairs.  There are, of course, many ways to represent sequences in
terms of pairs.  One particularly straightforward representation is illustrated
in \link{Figure 2.4}, where the sequence 1, 2, 3, 4 is represented as a chain
of pairs.  This structure is sometimes called a \newterm{linked list}, denoting
the idea that each element is linked to the next.  The \code{car} of each pair
is the corresponding item in the chain, and the \code{cdr} of the pair is the
next pair in the chain.  The \code{cdr} of the final pair signals the end of
the sequence by pointing to a distinguished value that is not a pair,
represented in box-and-pointer diagrams as a diagonal line and in programs as
the value of the variable \code{nil}.  The entire sequence is constructed by
nested \code{cons} operations:

\begin{codeBlock}
(cons 1
      (cons 2
            (cons 3
                  (cons 4 nil))))
\end{codeBlock}

% figure 2.4
\begin{imageFigure}{The sequence 1, 2, 3, 4 represented as a chain of pairs.}
  \includegraphics[width=76mm]{assets/figures/chapter_2/figure_2.4c.pdf}
\end{imageFigure}

Such a sequence of pairs, formed by nested \code{cons}es, is called a
\newterm{list}, and Scheme provides a primitive called \code{list} to help in
constructing lists.  The above sequence could be produced by
\code{(list 1 2 3 4)}.  In general,

\begin{syntaxForm}
(list \meta{a1} \meta{a2} ... \meta{an})
\end{syntaxForm}

\noindent
is equivalent to

\begin{syntaxForm}
(cons \meta{a1}
      (cons \meta{a2}
            (cons ...
                  (cons \meta{an}
                        nil)...)))
\end{syntaxForm}

\noindent
\begin{buildFromScarcity}
Scheme must think of a list this way because the pair is the only compound
structure it has.  A list is not a thing in its own right; it is a discipline
for arranging pairs, and the diagram in \link{Figure 2.4} is the whole of the
idea.

Python need not.  It supplies a sequence of its own as a core structure of the
language---the \newterm{tuple}---and we have already met it: the pair we built
in \link{Section 2.1} was a tuple that happened to have two elements.  A tuple
is written as a comma-separated series of objects, numbers in this case, and it
may have as many as we like.  So the sequence above is simply

\begin{codeBlock}
>>> (1, 2, 3, 4)
(1, 2, 3, 4)
\end{codeBlock}

\noindent
with no chain and no terminator, because there is nothing to terminate.  The
reader may wonder why we do not call this a list.  Python reserves that word:
its \code{list} is a different type, one whose contents can be changed after it
is made, and we shall take it up in \link{Chapter 3} when mutability is the
subject.  Until then a sequence for us is a tuple, and when this book says
``list'' in the original's sense, a tuple is what we shall build.

We shall keep the chain of pairs all the same.  It is the right picture to have
in mind when writing a procedure that works through a sequence a piece at a
time, which is nearly every procedure in this section---a first element and a
rest, over and over, until there is no rest left.  We simply decline to build
it.
\end{buildFromScarcity}

Where Scheme provides a built-in procedure \code{list} for assembling a
sequence, Python asks only for commas.  There is no constructor to call: string
the elements together with commas and the sequence is made, so that the means
of combination is a piece of punctuation rather than a procedure.

Python then prints a tuple as one writes it---the elements separated by commas,
enclosed in parentheses---so that what comes back on the screen is an
expression one could type in again.  Thus the data object of \link{Figure 2.4}
is printed as \code{(1, 2, 3, 4)}:

\begin{codeBlock}
>>> one_through_four = (1, 2, 3, 4)
>>> one_through_four
(1, 2, 3, 4)
\end{codeBlock}

\noindent
We can think of \code{[0]} as selecting the first item in the sequence, and of
\code{[1:]} as selecting the subsequence consisting of all but the first item.
The first of these we have seen already: it is how \code{numer} reached into a
rational number in \link{Section 2.1.1}.  The second is new.  Written between
brackets, a colon asks for a \newterm{slice}---a run of the sequence rather
than a single element---and \code{[1:]} means ``from element 1 to the end,''
which is to say everything after the first.  Repeated selections extract the
second, third, and subsequent items.\footnote{Nested applications of \code{car}
and \code{cdr} are cumbersome enough to write that Lisp dialects provide
abbreviations for them---\code{(cadr \meta{arg})} for \code{(car (cdr
\meta{arg}))}, and so on, each \code{a} standing for a \code{car} and each
\code{d} for a \code{cdr}.  The names persist because combinations like
\code{cadr} are pronounceable.  Python needs no such contraction, since
\code{items[1]} says directly what \code{(cadr items)} says by composition.}
To make a sequence like the original one but with an additional item at the
beginning, we join a one-element tuple to it with \code{+}.

\begin{codeBlock}
>>> one_through_four[0]
1
>>> one_through_four[1:]
(2, 3, 4)
>>> one_through_four[1:][0]
2
>>> (10,) + one_through_four
(10, 1, 2, 3, 4)
>>> (5,) + one_through_four
(5, 1, 2, 3, 4)
\end{codeBlock}

\noindent
Note the comma in \code{(10,)}.  As we saw in \link{Section 2.1.1}, it is the
comma and not the parentheses that makes a tuple, so \code{(10)} would be
merely the number 10 in brackets and \code{(10) + one\_through\_four} would be
an error rather than a longer sequence.

\begin{notTaste}
Throughout this chapter we shall add new elements at the \emph{front} of a
sequence, as \code{(10,) + one\_through\_four} does, because that is what the
original's algorithms do and following them is the point.  It should be said
plainly that this is not how Python is ordinarily written.  A Python programmer
building up a sequence adds to the end, and there is a reason beyond habit: a
pair in Scheme can point at an existing list and so \code{cons} costs nothing,
whereas a tuple cannot share, and joining one to another copies both.  We will
return to what that costs us in \link{Section 2.2.3}, when we have enough
procedures written for the cost to be worth measuring.
\end{notTaste}

The value of \code{nil}, used to terminate the chain of pairs in Scheme, can be
thought
of as a sequence of no elements, the \newterm{empty list}.  For us it is the
empty tuple, written \code{()}---a pair of parentheses with nothing between
them, and the one tuple that needs no comma.\footnote{It is remarkable how much
energy in the standardization of Lisp dialects was dissipated in arguments that
are literally over nothing: should \code{nil} be an ordinary name?  Should its
value be a symbol?  A list?  A pair?  The authors of the original, who had
endured too many language standardization brawls, preferred to avoid the entire
issue, and from \link{Section 2.3} onward denote the empty list \code{'()}.
Python was spared the argument.  \code{()} is a tuple of no elements, and there
is nothing further to say about it.}

\subsubsection*{List operations}

The use of sequences of elements is accompanied by conventional programming
techniques for manipulating them by successively taking the rest---what the
original calls ``\code{cdr}ing down'' a list.  For example, the procedure
\code{list\_ref} takes as arguments a sequence and a number \( n \) and returns
the \( n^{\mathrm{th}} \) item of the sequence.  It is customary to number the
elements beginning with 0.  The method for computing \code{list\_ref} is the
following:

\begin{itemize}
\item
For \( n = 0 \), \code{list\_ref} should return the first element.

\item
Otherwise, \code{list\_ref} should return the \( (n - 1) \)-st item of the rest
of the sequence.
\end{itemize}

\begin{codeBlock}
>>> @tail_recursive
... def list_ref(items, n):
...     if n == 0:
...         return items[0]
...     else:
...         return TailCall(list_ref, items[1:], n - 1)

>>> squares = (1, 4, 9, 16, 25)
>>> list_ref(squares, 3)
16
\end{codeBlock}

\noindent
\begin{revealTheSurface}
That procedure walks the sequence one element at a time, and Python has a
statement for doing exactly that:

\begin{codeBlock}
>>> def list_ref(items, n):
...     for index, item in enumerate(items):
...         if index == n:
...             return item
\end{codeBlock}

\noindent
The \code{for} statement takes each element of a sequence in turn and runs its
body with the name bound to that element.  Written plainly it would be
\code{for item in items:}, which gives us the elements but not their positions;
\code{enumerate} supplies both, handing back the index alongside the element at
each step, so that we can tell when we have arrived at the one we want.

Python has the whole procedure already, of course.  Selecting an element by
position is what the brackets are for, and \code{list\_ref} is spelled

\begin{codeBlock}
>>> squares[3]
16
\end{codeBlock}

\noindent
This is the same notation we used above for \code{[0]}, which was never a
special case---\code{[0]} is simply \code{list\_ref} with \( n \) fixed at
zero.  Nor is the slice \code{[1:]} anything new: it is the same brackets asked
for a run instead of a single element.
\end{revealTheSurface}

Often we work down the whole sequence, and must be able to tell when nothing is
left of it.  To aid in this, Scheme provides a primitive predicate
\code{null?}, which tests whether its argument is the empty list.  Python
provides no such procedure, and needs none.  We can compare against the empty
tuple directly and write \code{items == ()}; but an empty tuple is itself
false, so the ordinary way to ask is \code{not items}.\footnote{An empty tuple,
an empty string, an empty dictionary and the number zero are all false in
Python, and their non-empty counterparts are true, as is every other number.
This is worth knowing well beyond sequences: it is why so many Python
conditionals test a value directly rather than comparing it with anything.} 
The procedure \code{length}, which returns the number of items in a sequence,
illustrates this typical pattern of use:

\begin{codeBlock}
>>> def length(items):
...     if not items:
...         return 0
...     else:
...         return 1 + length(items[1:])

>>> odds = (1, 3, 5, 7)
>>> length(odds)
4
\end{codeBlock}

\noindent
The \code{length} procedure implements a simple recursive plan.  The reduction
step is:

\begin{itemize}
\item
The \code{length} of any sequence is 1 plus the \code{length} of the rest of
the sequence.
\end{itemize}

\noindent
This is applied successively until we reach the base case:

\begin{itemize}
\item
The \code{length} of the empty sequence is 0.
\end{itemize}

\noindent
We could also compute \code{length} in an iterative style:

\begin{codeBlock}
>>> def length(items):
...     @tail_recursive
...     def length_iter(a, count):
...         if not a:
...             return count
...         else:
...             return TailCall(length_iter, a[1:], count + 1)
...     return length_iter(items, 0)
\end{codeBlock}

\noindent
\begin{revealTheSurface}
Or with a \code{for} statement.  Here we want only to count the elements and
never to look at one, so there is no need to name them:

\begin{codeBlock}
>>> def length(items):
...     count = 0
...     for _ in items:
...         count = count + 1
...     return count
\end{codeBlock}

\noindent
The underscore is an ordinary name, but by long convention it is the name a
Python programmer uses for a value that must be received and will not be used.
Reading it, one knows at once that the body does not care what the elements
are.

And Python counts sequences already, with the built-in procedure \code{len}:

\begin{codeBlock}
>>> len(odds)
4
\end{codeBlock}
\end{revealTheSurface}

Another conventional programming technique is to build up an answer while
working down a sequence, as in the procedure \code{append}, which takes two
sequences as arguments and combines their elements to make a new one:

\begin{codeBlock}
>>> append(squares, odds)
(1, 4, 9, 16, 25, 1, 3, 5, 7)
>>> append(odds, squares)
(1, 3, 5, 7, 1, 4, 9, 16, 25)
\end{codeBlock}

\noindent
\code{append} is also implemented using a recursive plan.  To \code{append}
sequences \code{list1} and \code{list2}, do the following:

\begin{itemize}
\item
If \code{list1} is the empty sequence, then the result is just \code{list2}.

\item
Otherwise, \code{append} the rest of \code{list1} and \code{list2}, and join
the first element of \code{list1} onto the front of the result:
\end{itemize}

\begin{codeBlock}
>>> def append(list1, list2):
...     if not list1:
...         return list2
...     else:
...         return (list1[0],) + append(list1[1:], list2)
\end{codeBlock}

\noindent
\begin{revealTheSurface}
The \code{+} that joins a one-element tuple to the front of another is not
limited to one element.  It joins any two tuples, which is the whole of
\code{append}:

\begin{codeBlock}
>>> squares + odds
(1, 4, 9, 16, 25, 1, 3, 5, 7)
\end{codeBlock}
\end{revealTheSurface}

\begin{exerciseGroup}
\phantomsection\label{Exercise 2.17}
\exerciseHeading{Exercise 2.17:} Define a procedure
\code{last\_pair} that returns the sequence containing only the last element of
a given (nonempty) sequence:

\begin{codeBlock}
>>> last_pair((23, 72, 149, 34))
(34,)
\end{codeBlock}

Write it twice, once as a recursive procedure and once with a \code{for}
statement.  Then observe that Python will do it for you: a negative index
counts from the end, so \code{items[-1]} is the last \emph{element}, \code{34},
and the slice \code{items[-1:]} is the one-element sequence \code{(34,)} that
this exercise asks for.  Which of the three would you rather read in six
months?

\phantomsection\label{Exercise 2.18}
\exerciseHeading{Exercise 2.18:} Define a procedure \code{reverse}
that takes a sequence as argument and returns a sequence of the same elements
in reverse order:

\begin{codeBlock}
>>> reverse((1, 4, 9, 16, 25))
(25, 16, 9, 4, 1)
\end{codeBlock}

Write it as a recursive procedure.  Python has this one too, in two spellings:
\code{tuple(reversed(items))}, and the slice \code{items[::-1]}, whose third
number is a step and whose \( -1 \) means ``backwards.''

\phantomsection\label{Exercise 2.19}
\exerciseHeading{Exercise 2.19:} Consider the change-counting
program of \link{Section 1.2.2}.  It would be nice to be able to easily change
the currency used by the program, so that we could compute the number of ways
to change a British pound, for example.  As the program is written, the
knowledge of the currency is distributed partly into the procedure
\code{first\_denomination} and partly into the procedure \code{count\_change}
(which knows that there are five kinds of U.S. coins).  It would be nicer to be
able to supply a sequence of coins to be used for making change.

We want to rewrite the procedure \code{cc} so that its second argument is a
sequence of the values of the coins to use rather than an integer specifying
which coins to use.  We could then have sequences that defined each kind of
currency:

\begin{codeBlock}
>>> us_coins = (50, 25, 10, 5, 1)
>>> uk_coins = (100, 50, 20, 10, 5, 2, 1, 0.5)
\end{codeBlock}

We could then call \code{cc} as follows:

\begin{codeBlock}
>>> cc(100, us_coins)
292
\end{codeBlock}

To do this will require changing the program \code{cc} somewhat.  It will still
have the same form, but it will access its second argument differently, as
follows:

\begin{codeBlock}
>>> def cc(amount, coin_values):
...     if amount == 0:
...         return 1
...     elif amount < 0 or is_empty(coin_values):
...         return 0
...     else:
...         return (cc(amount, except_first_denomination(coin_values))
...                 + cc(amount - first_denomination(coin_values),
...                      coin_values))
\end{codeBlock}

Define the procedures \code{first\_denomination},
\code{except\_first\_denomination}, and \code{is\_empty} in terms of primitive
operations on sequences.\footnote{The original calls the last of these
\code{no-more?}.  Scheme marks a predicate with a trailing question mark, which
Python cannot spell; \code{is\_no\_more} is the literal rendering and reads
badly, so we have used \code{is\_empty}, which says the same thing.}  Does the
order of the sequence \code{coin\_values} affect the answer produced by
\code{cc}?  Why or why not?

\phantomsection\label{Exercise 2.20}
\exerciseHeading{Exercise 2.20:} The procedures \code{+},
\code{*}, and \code{list} take arbitrary numbers of arguments.  One way to
define such procedures in the original is to use \newterm{dotted-tail
notation}: a parameter list with a dot before the last parameter name
indicates that the initial parameters (if any) will have as values the initial
arguments, as usual, but the final parameter's value will be a list of any
remaining arguments.

\begin{syntaxForm}
(define (f x y . z) \meta{body})
(define (g . w) \meta{body})
\end{syntaxForm}

\noindent
Python writes the dot as a star, and puts it on the parameter rather than
before it:

\begin{codeBlock}
>>> def f(x, y, *z):
...     return (x, y, z)

>>> f(1, 2, 3, 4, 5, 6)
(1, 2, (3, 4, 5, 6))
\end{codeBlock}

\noindent
Called with two or more arguments, \code{f} binds \code{x} to 1, \code{y} to 2,
and \code{z} to the tuple \code{(3, 4, 5, 6)}.  Note that inside the body
\code{z} is an ordinary tuple, and the star does not appear again: it marks the
parameter when the procedure is defined, and thereafter \code{z} is a sequence
like any other.  A procedure whose whole parameter list is starred takes zero
or more arguments:

\begin{codeBlock}
>>> def g(*w):
...     return w

>>> g(1, 2, 3, 4, 5, 6)
(1, 2, 3, 4, 5, 6)
>>> g()
()
\end{codeBlock}

\noindent
Such a parameter is conventionally named \code{args}, and one reads
\code{*args} in Python programs constantly.  \code{w} and \code{z} are just
names, as \code{x} and \code{y} are.

Use this notation to write a procedure \code{same\_parity} that takes one or
more integers and returns a sequence of all the arguments that have the same
even-odd parity as the first argument.  For example,

\begin{codeBlock}
>>> same_parity(1, 2, 3, 4, 5, 6, 7)
(1, 3, 5, 7)
>>> same_parity(2, 3, 4, 5, 6, 7)
(2, 4, 6)
\end{codeBlock}
\end{exerciseGroup}

\subsubsection*{Mapping over lists}

One extremely useful operation is to apply some transformation to each element
in a sequence and generate the sequence of results.  For instance, the
following procedure scales each number in a sequence by a given factor:

\begin{codeBlock}
>>> def scale_list(items, factor):
...     if not items:
...         return ()
...     else:
...         return ((items[0] * factor,)
...                 + scale_list(items[1:], factor))

>>> scale_list((1, 2, 3, 4, 5), 10)
(10, 20, 30, 40, 50)
\end{codeBlock}

\noindent
We can abstract this general idea and capture it as a common pattern expressed
as a higher-order procedure, just as in \link{Section 1.3}.  The higher-order
procedure here is called \code{map}.  \code{map} takes as arguments a procedure
of one argument and a sequence, and returns a sequence of the results produced
by applying the procedure to each element:

\begin{codeBlock}
>>> def map(procedure, items):
...     if not items:
...         return ()
...     else:
...         return ((procedure(items[0]),)
...                 + map(procedure, items[1:]))

>>> map(abs, (-10, 2.5, -11.6, 17))
(10, 2.5, 11.6, 17)
>>> map(lambda x: x * x, (1, 2, 3, 4))
(1, 4, 9, 16)
\end{codeBlock}

\noindent
\begin{revealTheSurface}
\code{map} is a built-in procedure in Python, and the definition above shadows
it---which is why we were able to write it at all.  To see the built-in one
again we must put our own away, which \code{del} does:

\begin{codeBlock}
>>> del map
\end{codeBlock}

\noindent
The built-in \code{map} differs from ours in two ways.  The first is that it is
more general, as the original's is: given a procedure of \( n \) arguments and
\( n \) sequences, it applies the procedure to all the first elements, then all
the second elements, and so on.

The second is that it does not return a sequence:

\begin{codeBlock}
>>> map(abs, (-10, 2.5, -11.6, 17))
<map object at 0x102a4c1f0>
\end{codeBlock}

\noindent
What comes back is an object that has not computed anything yet and will
produce each result only when asked for it---an \newterm{iterator}.  This is
called being \newterm{lazy}, and it is why nothing was printed: no one has
asked.  To get a sequence, ask for one, by handing the iterator to
\code{tuple}:

\begin{codeBlock}
>>> tuple(map(abs, (-10, 2.5, -11.6, 17)))
(10, 2.5, 11.6, 17)
>>> tuple(map(lambda x: x * x, (1, 2, 3, 4)))
(1, 4, 9, 16)
>>> tuple(map(lambda x, y: x + 2 * y, (1, 2, 3), (4, 5, 6)))
(9, 12, 15)
\end{codeBlock}

\noindent
Laziness is a real advantage---a map over a very long sequence need never build
the whole result---and a real trap, which we shall spring deliberately in
\link{Exercise 2.23}.
\end{revealTheSurface}

Now we can give a new definition of \code{scale\_list} in terms of \code{map}:

\begin{codeBlock}
>>> def scale_list(items, factor):
...     return tuple(map(lambda x: x * factor, items))
\end{codeBlock}

\noindent
\code{map} is an important construct, not only because it captures a common
pattern, but because it establishes a higher level of abstraction in dealing
with sequences.  In the original definition of \code{scale\_list}, the
recursive structure of the program draws attention to the element-by-element
processing of the sequence.  Defining \code{scale\_list} in terms of \code{map}
suppresses that level of detail and emphasizes that scaling transforms a
sequence of elements to a sequence of results.  The difference between the two
definitions is not that the computer is performing a different process (it
isn't) but that we think about the process differently.  In effect, \code{map}
helps establish an abstraction barrier that isolates the implementation of
procedures that transform sequences from the details of how the elements are
extracted and combined.  Like the barriers shown in \link{Figure 2.1}, this
abstraction gives us the flexibility to change the low-level details of how
sequences are implemented, while preserving the conceptual framework of
operations that transform sequences to sequences.  \link{Section 2.2.3} expands
on this use of sequences as a framework for organizing programs.

\begin{exerciseGroup}
\phantomsection\label{Exercise 2.21}
\exerciseHeading{Exercise 2.21:} The procedure \code{square\_list}
takes a sequence of numbers as argument and returns a sequence of the squares
of those numbers.

\begin{codeBlock}
>>> square_list((1, 2, 3, 4))
(1, 4, 9, 16)
\end{codeBlock}

Here are two different definitions of \code{square\_list}.  Complete both of
them by filling in the missing expressions:

\begin{syntaxForm}
def square_list(items):
    if not items:
        return ()
    else:
        return \meta{??} + \meta{??}

def square_list(items):
    return tuple(map(\meta{??}, \meta{??}))
\end{syntaxForm}

\phantomsection\label{Exercise 2.22}
\exerciseHeading{Exercise 2.22:} Louis Reasoner tries to rewrite
the first \code{square\_list} procedure of \link{Exercise 2.21} so that it
evolves an iterative process:

\begin{codeBlock}
>>> def square_list(items):
...     answer = ()
...     for item in items:
...         answer = (square(item),) + answer
...     return answer
\end{codeBlock}

Unfortunately, defining \code{square\_list} this way produces the answer in the
reverse order of the one desired.  Why?

Louis then tries to fix his bug by interchanging the arguments to \code{+}:

\begin{codeBlock}
>>> def square_list(items):
...     answer = ()
...     for item in items:
...         answer = answer + (square(item),)
...     return answer
\end{codeBlock}

This time it works.  Why?

The original poses this exercise in Scheme, where the second version does
\emph{not} work: interchanging the arguments to \code{cons} does not append to
a list, it builds a pair whose first element is the accumulated list and whose
second is a number, and the result is not a list at all.  Explain why the same
interchange succeeds here.  What is it about \code{+} on tuples that
\code{cons} on pairs does not do?

Then count the work.  Each version joins two tuples once per element, and
joining two tuples copies both.  How many element-copies does each version
perform on a sequence of \( n \) elements?  How many would the original's
\code{cons} perform?  Keep the answer in mind for \link{Section 2.2.3}.

\phantomsection\label{Exercise 2.23}
\exerciseHeading{Exercise 2.23:} The procedure \code{for\_each} is
similar to \code{map}.  It takes as arguments a procedure and a sequence of
elements.  However, rather than forming a sequence of the results,
\code{for\_each} just applies the procedure to each of the elements in turn,
from left to right. The values returned by applying the procedure to the
elements are not used at all---\code{for\_each} is used with procedures that
perform an action, such as printing.  For example,

\begin{codeBlock}
>>> for_each(print, (57, 321, 88))
57
321
88
\end{codeBlock}

Give a recursive implementation of \code{for\_each}.

Then consider that Python appears to have made the procedure unnecessary, since
the \code{for} statement does precisely this:

\begin{codeBlock}
>>> for item in (57, 321, 88):
...     print(item)
57
321
88
\end{codeBlock}

Now try it with \code{map}, whose whole business is applying a procedure to
every element:

\begin{codeBlock}
>>> map(print, (57, 321, 88))
<map object at 0x1027ecd80>
\end{codeBlock}

Nothing is printed.  Explain, and say what \code{tuple(map(print, (57, 321,
88)))} would print and what it would return.  The original notes that the value
returned by \code{for\_each} can be something arbitrary; what does the
\code{for} statement return?  Between them these two questions explain why
Python keeps \code{map} and \code{for} as separate things, where the original's
\code{map} and \code{for-each} differ only in whether anyone keeps the results.
\end{exerciseGroup}
